\section{Introduction}\label{sec:intro}
Segmentation of sequential data, also known as change-point detection, is a fundamental problem in the field
of unsupervised learning, and has applications in diverse fields such
as speech processing \cite{brent1999speech,phonemeSeg:2008,Shriberg2000127},
text processing \cite{beeferman1999statistical}, bioinformatics \cite{olshen2004circular}
and network anomaly detection \cite{levy2009detection}, to name a few.
We are interested in formulating the segmentation task as a convex optimization
problem that avoids issues such as local-minima or sensitivity to
initializations. In addition, we want to explicitly incorporate robustness to outliers.
% These issues are often encountered when using other
%approaches, such as maximization of the likelihood of the data given
%some generative model.
Given a sequence of samples $\{\vxii\}_{i=1}^n$, for
$\vxii\in\mathbb{R}^d$, our goal is to segment it into a few
subsequences, where each subsequence is homogeneous under some criterion. Our starting point is a convex objective
that minimizes the sum of squared distances of samples $\vxii$ from each sample's associated `centroid`, $\vmuii$. Identical adjacent $\vmuii$s identify their corresponding samples as belonging to the same segment.
In addition, some of the samples are allowed to be identified as outliers, allowing reduced loss on these samples. Two regularization terms are added to the objective, in order to constrain the number of detected segments and outliers, respectively. %These regularization terms can be seen as a convex relaxation of counting the number of segments and outliers.

%that was considered by
%\cite{vert2010fast}, in which each sample $\vxii$ is associated with a
%representative vector $\vmui{i}$. This
%objective captures two completing properties. First, that $\vmuii$
%will be close to $\vxii$ in the Euclidean sense, and second, that many
%adjacent $\vmuii$'s will be identical, a property which identifies
%them as belonging to the same segment, while controlling the actual number of segments $K$.
We propose two algorithms
based on this formulation, both alternate between detecting outliers, which is solved analytically, and solving the problem
with modified samples %and no outliers
, which can be solved iteratively.
The first algorithm, denoted by Outlier-Robust Convex Sequential (ORCS)
segmentation, solves the optimization problem exactly, while the
second is a top-down hierarchical version of the algorithm, called
TD-ORCS. We also derive a weighted version of this algorithm, denoted by WTD-ORCS. We show that for the case of $K=2$ segments and no outliers, a specific choice of the weights leads to a solution which recovers the exact solution of an un-relaxed optimization problem.

%We first propose an analytical solution to the optimization problem for the case of $K=2$. We analyze this solution and derive a bound on the probability that it is far from the correct boundary. We use this bound in order to formulate a weighted version of this analytical solution, and we show that for a specific choice of the weights, the solution recovers the exact solution of the un-relaxed optimization problem. We then propose an approximated, top-down greedy algorithm for a general $K$.

%We discuss the solution to the optimization problem, and why standard proximal gradient methods like the ones presented by ~\cite{bach2011convex,Beck_FISTA:2009} may not be feasible. We propose two ways for addressing this issue. One via the dual problem, and the other via hierarchical top-down greedy approach.  We derive an exact solution for the simple case of $K=2$ segments. We also analyze this solution and show that if indeed $K=2$ and data was generated by a specific stochastic process, with high-probability, the segmentation induced from the optimal solution will not be too far from the correct one. We further use the bound derived in order to formulate a weighted version of the top-down algorithm.

%As shown in the experiments below,
%the performance of these algorithms %derived from this formulation
%degrades fast and severely in
%the context of non-stationary (outliers) noise. We thus derive a new convex
%objective that takes outliers into consideration, and introduce an
%additional set of variables $\{\vzii\}$ which when non-zero,
%identifies a data sample as an outlier.
%We use a relaxation for
%counting the number of such samples, which allows for decreasing the
%squared loss suffered when a given data sample $\vxii$ is indeed an
%outlier.
%The number of allowed outliers is constrained to avoid
%identifying all data samples as outliers. We propose two algorithms
%based on this formulation, both alternate between solving the problem
%with modified samples and no outliers, and detecting outliers. The
%first algorithm, denoted by Outlier-Robust Convex Sequential (ORCS)
%segmentation, solves the optimization problem exactly, while the
%second is a top-down hierarchical version of the algorithm, called
%TD-ORCS. We also derive a weighted version of this algorithm, denoted by WTD-ORCS.
%The hierarchical algorithm is greedy and does not solve the problem
%exactly, yet it yields better empirical performance.

We evaluate the
performance of the proposed algorithms on %synthetic data, and on
two speech segmentation tasks,
%One is segmentation of fundamental speech units called phonemes, and the second is segmentation of a
%radio talk-show recording. The evaluation is done
for both clean sources and sources
contaminated with added non-stationary noise. %manually added outliers.
Our algorithms outperform other algorithms in both the clean and outlier-contaminated setting.
Finally, based on the empirical results, we propose a heuristic approach for approximating the number of outliers.
%Our algorithms outperformed other algorithms when more than
%$10\%$ of the segments were contaminated. Furthermore, TD-ORCS
%performance is as high as any algorithms designed for the no-outliers
%case, and as high as any algorithm designed for the outliers case.

%\paragraph*{Notation}
{\bf Notation}: The samples to be segmented are
denoted by $\{\vxii\in\mathbb{R}^d\}_{i=1}^n$, and their associated
quantities (both variables and solutions) are
$\vmuii,\vzii\in\mathbb{R}^d$. The same notation with no subscript,
$\vmu$, denotes the collection of all $\vmuii$'s.
%, or an arbitrary
%member of this collection, where it is clear from context which
%meaning is used.
The same holds for $\vx,\vz$. We abuse notation and
use $\vmuii$ to refer to both the `centroid' vector of a segment (these
are not center of mass, due to the regularization term), and to the
indexes of measurements assigned to that segment.
%Formally, given some
%vector $\vmuii$ we denote by  $S(\vmuii)$ the maximal set of all consecutive vectors
%equal to it. %, $S(\vmuii)=\{ j ~:~ \vmui{j}=\vmui{i} \}$.
%This set
%satisfies that if $j<k \in S(\vmuii)$ then also $j+1 ... k-1 \in
%S(\vmuii)$, and $\vmui{ \min_j \in S(\vmuii) - 1} \neq \vmuii$ and
%$\vmui{ \max_j \in S(\vmuii) + 1} \neq \vmuii$. We will abuse notation
%below and identify $\vmuii$ with $S(\vmuii)$.
